% !TEX TS-program = lualatex
Следующим этапом развития поисковых алгоритмов в играх стало использование вероятностных и стохастических методов, таких как Monte Carlo Tree Search (MCTS), которые получили широкое распространение в задачах с высокой степенью неопределённости и большой размерностью пространства состояний. Одной из ключевых особенностей данного алгоритма является использование памяти: его ход организован таким образом, чтобы с помощью специально введённых параметров, таких как количество сыгранных партий и число выигранных матчей, анализировать ветви дерева игры \cite{monte-carlo-tree}.

Шаги \refalg{mcts-algorithm} можно представить следующим образом.

\begin{alg}{Алгоритм построения дерева MCTS}{mcts-algorithm}
    \While{игра не завершена}
        \State Выбрать ход соперника (если он не существует в дереве --- добавить его)
        \State Расширить дерево: добавить ход агента с <<нулевым результатом>>
        \State Провести симуляцию: доиграть партию до конца из текущего состояния
        \State Распространить результат симуляции от текущего узла к корню
    \EndWhile
\end{alg}

На \refimage{mcts-example} cхематически представлены шаги ключевых этапов алгоритма. Процесс начинается с этапа выбора, где алгоритм проходит от корня дерева к наиболее перспективным узлам на основе определённой стратегии; выбранные узлы выделены тёмным цветом и содержат соотношение побед к посещениям (например, 7/10). Далее следует расширение, на котором к дереву добавляется новый узел, представляющий ранее не исследованное действие или состояние --- на изображении это узел со значением 0/0. Затем происходит симуляция --- случайное проигрывание игры или сценария от нового узла до финала, чтобы смоделировать возможный результат, например, поражение (0/1). Наконец, на этапе обратного распространения результат симуляции передаётся вверх по дереву, обновляя статистику всех узлов на этом пути (например, узел изменяется с 7/10 на 7/11), что позволяет улучшить качество принятия решений в последующих итерациях.

\image[width=0.7\textwidth]{Шаги алгоритма <<Поиск по дереву Монте-Карло>>}{mcts-example}{mcts-example}

